% Prologue. Blake's voice at the seams. Target ~600 words, 2-4 pages.
%
% Brief from the plan:
%   - Why these letters survived to you.
%   - What it meant to find them.
%   - Joan and Gene as people, briefly. (Where they met, what kind of people
%     they were, how long they lived after these letters end.)
%   - The condition you found the letters in.
%   - The frame you want the reader to hold while they read.
%
% Style guidance: present at the seams, not woven throughout. Do not
% summarize the war here; leave Pearl Harbor and Tassafaronga to be lived
% inside the body of the book. The prologue earns the reader's trust and
% then steps back.

\cleardoublepage
\thispagestyle{laempty}
\vspace*{0.10\textheight}
\begin{center}
    \eyebrow{Prologue}\par
    \vspace{0.6em}
    \brassrule[36pt]\par
\end{center}
\vspace{2em}

\noindent
% =====================================================================
% TODO[blake]: REPLACE THIS PARAGRAPH WITH YOUR PROLOGUE.
% =====================================================================
%
% Suggested opening: a single sentence about how the letters came to you.
% Then 4--6 paragraphs in the order above. Closes with the frame you want
% the reader to carry into the book.
% ---------------------------------------------------------------------

\lettrine[lines=3, slope=0pt, findent=-0.2em, nindent=0pt]
{\displayfont\itshape T}{he letters in this book}
came to me in a zip file in 2022, scanned by a relative who knew
they should not be left where the next flood or the next move could
take them. There were eighty-six photographs, a typed transcription
by another relative, and a handwritten note that said, simply,
\emph{Gene and Joan}. \textsc{[\,placeholder\,]}

\vspace{0.6em}

\noindent
\textsc{[\,placeholder \,---\, two to three more paragraphs in Blake's voice. Aim for the
reader to know who Joan was, who Gene was, where they lived, how
they met, and that they were married for forty-nine years after
the war ended, without yet revealing the moments that the body of
the book is going to live through.\,]}

\vspace{0.6em}

\noindent
\textsc{[\,placeholder \,---\, the closing of the prologue. Hand the reader off to
April 14, 1940. One thought or one image. Sign it.\,]}

\vspace{2em}
\begin{flushright}
    {\handfont\Large Blake Morris}\par
    \vspace{0.2em}
    {\monofont\footnotesize\color{walnutmute}\addfontfeature{LetterSpace=22}
        SOMERSET, KENTUCKY \, \textbf{\textperiodcentered} \, \the\year}
\end{flushright}
\vfill
